\section*{Chapitre \ref{ch:b2:oceanic-islands} --- Îles océaniques}

\begin{solution}{exo:b2:oceanic-islands:1}
(a) Stade pré-bouclier, comme Loihi aujourd'hui. (b) \omterm{def:b2:oceanic-islands:stages}{Stade bouclier}. (c) Stade post-bouclier. (d) Stade de
rajeunissement. (e) Plus aucun stade éruptif~: le volcan s'est enfoncé sous son récif, qui est devenu un \omterm{def:b2:oceanic-islands:reefs}{atoll}.
\end{solution}

\begin{solution}{exo:b2:oceanic-islands:2}
(a) $2500 + 350\sqrt9 = \qty{3550}{m}$. (b) $350\times(\sqrt{25} - \sqrt9) = 350\times2 =
\qty{700}{m}$. (c) $350/(2\times5) = \qty{35}{m}$ par million d'années.
\end{solution}

\begin{solution}{exo:b2:oceanic-islands:3}
(a) $5000/17 = \qty{294}{ka}$ autour des îles principales~; $30\,000/70 = \qty{429}{ka}$
le long de la dorsale. (b) Des glissements échappent à la détection~: un volcan plus
jeune peut enfouir la cicatrice et les débris d'un plus ancien, et d'anciens dépôts sont
recouverts de sédiments. Le nombre réel est plus grand et l'intervalle réel plus court.
\end{solution}

\begin{solution}{exo:b2:oceanic-islands:4}
(a) $4158\times0.3048 = \qty{1267}{m}$~; $4610\times0.3048 = \qty{1405}{m}$~; $150\times0.3048 =
\qty{46}{m}$. (b) Le premier récif n'a pas poussé plus profond que \qty{46}{m}, et sa base se trouve aujourd'hui à
\qty{1267}{m}~: le volcan s'est enfoncé d'au moins $1267 - 46 = \qty{1221}{m}$ (davantage si le niveau
de la mer était plus haut alors). (c) L'enfoncement est au plus d'environ \qty{1267}{m}, sur au moins
\qty{33.9}{Ma}~: au plus $1267/33.9 = \qty{37}{m}$ par million d'années, soit \qty{0.037}{mm} par an.
\end{solution}

\begin{solution}{exo:b2:oceanic-islands:5}
(a) $\dot S = 175/\sqrt t$~: \qty{87.5}{m.Ma^{-1}} (\qty{0.0875}{mm.an^{-1}}) à \qty{4}{Ma},
\qty{35}{m.Ma^{-1}} (\qty{0.035}{mm.an^{-1}}) à \qty{25}{Ma}, \qty{25}{m.Ma^{-1}}
(\qty{0.025}{mm.an^{-1}}) à \qty{49}{Ma}. (b) $4.2/0.0875 = 48$ fois. (c) $2.6/0.0875 = 30$
fois~: la phase de charge enfonce une île bien plus vite que le refroidissement d'une jeune plaque.
\end{solution}

\begin{solution}{exo:b2:oceanic-islands:6}
(a) $\dot S = (3200/62.8)\,e^{-100/62.8} = 50.96\times0.2035 = \qty{10.37}{m}$ par million d'années,
soit \qty{0.0104}{m} par millier d'années~: la valeur mesurée, \qty{0.01}{m} par millier d'années. (b) $e^{-t/62.8}
= 10/50.96$, $t = 62.8\ln5.096 = \qty{102.3}{Ma}$. (c) Loi en racine~: $2500 + 350\times5 =
\qty{4250}{m}$~; loi exponentielle~: $6400 - 3200\times0.6716 = \qty{4251}{m}$~; l'écart est de
\qty{1}{m}.
\end{solution}

\begin{solution}{exo:b2:oceanic-islands:7}
(a) $\qty{8.6}{cm.an^{-1}} = \qty{86}{km.Ma^{-1}}$~; $86\times33 = \qty{2838}{km}$. (b) Oui~:
\qty{27.7}{Ma} est inférieur aux \qty{33}{Ma} au bout desquels les récifs se noient, donc son récif
doit être encore vivant. (c) $55\times1.609 = \qty{88.5}{km}$, soit $88.5/86 = \qty{1.0}{Ma}$ de
plus~: environ \qty{28.7}{Ma}. C'est un peu moins que \qty{33}{Ma}, mais le chiffre de \qty{33}{Ma}
est lui-même une valeur approchée pour toute la chaîne, et la croissance dépend aussi de la température
de l'eau à chaque \omterm{def:b2:oceanic-islands:reefs}{atoll}~: Kure, le dernier \omterm{def:b2:oceanic-islands:reefs}{atoll} vivant, est bien proche du \omterm{def:b2:oceanic-islands:darwin-point}{point de Darwin}.
\end{solution}

\begin{solution}{exo:b2:oceanic-islands:8}
(a) Boîte~: $30\times17\times2 = \qty{1020}{km^3}$~; pyramide, un tiers, \qty{340}{km^3}. (b) Un seul bloc représente à lui seul un dixième (boîte) ou un trentième (pyramide)
de l'ensemble du volcan Kilauea. (c) $23\,000\times
0.1 = \qty{2300}{km^3}$, soit un quart du Kilauea~; une estimation qui tient ou tombe avec
l'épaisseur supposée.
\end{solution}

\begin{solution}{exo:b2:oceanic-islands:9}
(a) $1000/2.6 = \qty{385}{ka}$. (b) Oui~: \qty{2.6}{m} par millier d'années font
\qty{2.6}{mm.an^{-1}}, moins que \qty{4.2}{mm.an^{-1}}, soit une marge de \qty{1.6}{mm.an^{-1}}.
(c) Remontée relative $12 + 2.6 = \qty{14.6}{mm.an^{-1}}$~; le récif prend du retard de $14.6 - 4.2 =
\qty{10.4}{mm.an^{-1}}$ pendant $16.5 - 8.2 = \qty{8.3}{ka}$~: son sommet finit à $10.4\times8.3 =
\qty{86}{m}$. (d) $(12 - 4.2)\times8.3 = \qty{65}{m}$. Dans les deux cas il reste au-dessus de la limite de croissance la plus profonde, mais
dans une eau nettement moins éclairée~: qu'il rattrape ou non dépend de la manière dont sa croissance
ralentit avec la profondeur.
\end{solution}

\begin{solution}{exo:b2:oceanic-islands:10}
(a) $S(t_1) = C(\sqrt{t_1 + \Delta} - \sqrt{t_1})$ a pour dérivée $\frac C2(1/\sqrt{t_1 +
\Delta} - 1/\sqrt{t_1}) < 0$~: elle décroît avec $t_1$. (b) $350\sqrt{20} = \qty{1565}{m}$~; $350\times
(\sqrt{70} - \sqrt{50}) = 350\times1.2955 = \qty{453}{m}$. (c) Non~: $\sqrt t$ est fini en $0$, donc
l'enfoncement sur un intervalle quelconque est fini, \qty{1565}{m} dans le cas (b). Seul le taux à l'instant $t = 0$
est infini, là où la loi n'est plus censée décrire le fond en détail.
\end{solution}

\begin{solution}{exo:b2:oceanic-islands:11}
(a) $30\times0.1 = \qty{3}{m}$. (b) Environ \qty{120}{m} si son sommet se trouvait au niveau de
la mer au haut niveau marin~: l'\omterm{def:b2:oceanic-islands:reefs}{atoll} devient une île calcaire, la pluie dissout l'aragonite des
coraux et la roche se recristallise en calcite, avec des cavités. (c) $120/3 = 40$ cycles, soit
\qty{4}{Ma}. (d) Les zones lessivées sont ces époques d'émersion, scellées sous le récif qui
a repoussé quand la mer est revenue. Sur de nombreux cycles, la mer revient à des niveaux hauts
comparables, donc $\Delta\eta$ est faible, le sommet du récif remonte au niveau de la mer à chaque
haut niveau, et l'épaisseur de la coiffe est fixée par l'enfoncement, comme l'énonce la proposition.
\end{solution}

\begin{solution}{exo:b2:oceanic-islands:12}
(a) $3200\times(e^{-50/62.8} - e^{-120/62.8}) = 3200\times(0.4510 - 0.1480) = \qty{970}{m}$.
(b) $150 + 970 = \qty{1120}{m}$. (c) $700\times1.8288 = \qty{1280}{m}$, soit \qty{160}{m}
de plus que l'estimation. (d) Le \omterm{def:b2:oceanic-islands:seamount}{guyot} a pu se noyer plus tôt, sur une plaque plus jeune,
et s'enfoncer davantage~; un volcan voisin a pu le charger~; son sommet a pu être rabaissé
par les vagues, ou être recouvert de sédiments~; la plaque située dessous peut avoir un
autre âge~; et le niveau de la mer à l'époque de la noyade n'était pas celui d'aujourd'hui.
\end{solution}

\begin{solution}{pb:b2:oceanic-islands:1}
\textbf{1.} $2500 + 350\sqrt{10} = 2500 + 350\times3.1623 = \qty{3607}{m}$.
\textbf{2.} $3607 + 1000 = \qty{4607}{m}$.
\textbf{3.} $V = \pi r^2h/3 = \pi\times50^2\times4.607/3 = \qty{12061}{km^3}$.
\textbf{4.} $0.85\times12\,061 = \qty{10252}{km^3}$ en \qty{1}{Ma}~: \qty{0.0103}{km^3} par an,
dix fois moins que le Kilauea aujourd'hui.
\textbf{5.} $95\,600/12\,061 = 7.9$~: environ huit.
\textbf{6.} Lave tholéiitique, issue du \omterm{def:b2:oceanic-islands:stages}{stade bouclier}~; les laves alcalines du stade post-bouclier
le coiffent.
\textbf{7.} $\qty{1}{km}/\qty{1}{Ma} = \qty{1}{m}$ par millier d'années, 2.6 fois plus lent qu'au
large de l'île d'Hawaï.
\textbf{8.} Oui~: $4.2 - 1.0 = \qty{3.2}{mm.an^{-1}}$ de marge.
\textbf{9.} $12 + 1 - 4.2 = \qty{8.8}{mm}$ par an.
\textbf{10.} $150/8.8 = \qty{17}{ka}$. Une terrasse récifale noyée sur le flanc, une marche du type
de celles observées au large de l'île d'Hawaï.
\textbf{11.} $350\times(\sqrt{22} - \sqrt{10}) = 350\times(4.6904 - 3.1623) = \qty{535}{m}$.
\textbf{12.} $1000 - 535 = \qty{465}{m}$, soit $465/12 = \qty{38.8}{m}$ par million d'années.
\textbf{13.} $350/(2\sqrt{10}) = \qty{55.3}{m.Ma^{-1}}$ au départ~; $350/(2\sqrt{22}) =
\qty{37.3}{m.Ma^{-1}}$ à la fin.
\textbf{14.} $2500 + 350\times4.6904 = \qty{4142}{m}$.
\textbf{15.} Les grandes îles hawaïennes s'élèvent jusqu'à environ quatre fois plus haut, donc il y a bien plus de
roche au-dessus du niveau de la mer à enlever~; leur \omterm{def:b2:oceanic-islands:stages}{stade bouclier} a aussi duré plus longtemps.
\textbf{16.} Les vagues, qui font reculer la côte et aplanissent une plate-forme d'abrasion marine~; les
glissements et les \omterm{def:b2:oceanic-islands:collapse}{effondrements de flanc}.
\textbf{17.} Un \omterm{def:b2:oceanic-islands:reefs}{récif-barrière}, avec un lagon entre le récif et le dernier rivage volcanique.
\textbf{18.} $350\times(\sqrt{70} - \sqrt{22}) = 350\times(8.3666 - 4.6904) = \qty{1287}{m}$.
\textbf{19.} $1287/48 = \qty{26.8}{m}$ par million d'années, soit \qty{0.027}{mm} par an~: $4.2/0.0268 =
157$ fois moins que ce que le récif pourrait croître.
\textbf{20.} \qty{20}{m} de plus que le calcaire au-dessus du basalte sous l'île Parry
(\qty{1267}{m}), et \qty{118}{m} de moins que sous Elugelab (\qty{1405}{m})~: de même ordre de grandeur.
\textbf{21.} $2500 + 350\times8.3666 = \qty{5428}{m}$.
\textbf{22.} $3200\times(e^{-70/62.8} - e^{-120/62.8}) = 3200\times(0.3280 - 0.1480) =
\qty{576}{m}$.
\textbf{23.} $150 + 576 = \qty{726}{m}$.
\textbf{24.} La coiffe récifale de l'\omterm{def:b2:oceanic-islands:reefs}{atoll} Petrel mesure \textbf{\qty{1287}{m}} d'épaisseur, soit environ \qty{1.3}{km}
de calcaire récifal, comme sous Enewetak.
\end{solution}
